\documentclass[conference]{IEEEtran}
\IEEEoverridecommandlockouts

\usepackage{cite}
\usepackage{amsmath,amssymb,amsfonts}
\usepackage{algorithmic}
\usepackage{graphicx}
\usepackage{textcomp}
\usepackage{xcolor}
\usepackage{listings}
\usepackage{booktabs}
\usepackage{url}
\usepackage[utf8]{inputenc}

\lstset{
  basicstyle=\ttfamily\footnotesize,
  breaklines=true,
  frame=single,
  language=[Sharp]C,
  keywordstyle=\color{blue},
  commentstyle=\color{green!60!black},
  stringstyle=\color{orange}
}

\def\BibTeX{{\rm B\kern-.05em{\sc i\kern-.025em b}\kern-.08em
    T\kern-.1667em\lower.7ex\hbox{E}\kern-.125emX}}

\begin{document}

\title{Implementimi i Algoritmit Post-Kuantik ML-KEM për Sigurinë Kriptografike në Epokën e Kompjuterave Kuantikë\\
{\large Implementation of Post-Quantum ML-KEM Algorithm for Cryptographic Security in the Quantum Computing Era}
}

\author{
\IEEEauthorblockN{Erand Kurtaliqi}
\IEEEauthorblockA{\textit{Programi Master, Siguria e Informacionit} \\
\textit{Universiteti i Prishtinës}\\
Prishtinë, Kosovë \\
erand.kurtaliqi@student.uni-pr.edu}
}

\maketitle

\begin{abstract}
Zhvillimi i kompjuterëve kuantikë paraqet një kërcënim thelbësor për sistemet aktuale kriptografike që bazohen në algoritme asimetrike si RSA dhe ECC. Ky punim prezanton implementimin praktik të algoritmit ML-KEM (Module-Lattice-Based Key-Encapsulation Mechanism), i standardizuar nga NIST si FIPS 203, për të siguruar komunikime të sigurta edhe në prani të kompjuterëve kuantikë. Sistemi i zhvilluar demonstron se kompjuterët klasikë mund të përdorin algoritme post-kuantike për të mbrojtur të dhënat nga sulmet e ardhshme kuantike. Implementimi përfshin një API në .NET 9 me BouncyCastle dhe një frontend Angular që ofron enkriptim hibrid ML-KEM + AES-GCM. Rezultatet tregojnë se ML-KEM ofron sigurinë e nevojshme me performancë të pranueshme për përdorim praktik.
\end{abstract}

\begin{IEEEkeywords}
Post-Quantum Cryptography, ML-KEM, CRYSTALS-Kyber, Lattice-Based Cryptography, NIST FIPS 203, AES-GCM, Key Encapsulation Mechanism
\end{IEEEkeywords}

\section{Hyrje}

Revolucioni i kompjuterave kuantikë po sjell ndryshime paradigmatike në fushën e sigurisë kibernetike. Algoritmi i Shor \cite{shor1994}, i publikuar në vitin 1994, demonstroi se kompjuterët kuantikë mund të thyejnë algoritmet e enkriptimit asimetrik si RSA dhe ECC në kohë polinomiale, duke i bërë të pavlefshme themelet e sigurisë së sotme dixhitale.

Sipas raporteve të fundit, kompjuterët kuantikë pritet të arrijnë aftësinë e thyrjes së RSA-2048 brenda dekadës së ardhshme \cite{nas2019}. Kjo krijon nevojën urgjente për migrimin në algoritme post-kuantike që janë rezistente ndaj sulmeve kuantike por që mund të ekzekutohen në kompjuterë klasikë.

NIST (National Institute of Standards and Technology) përfundoi procesin e standardizimit në gusht 2024, duke publikuar FIPS 203 \cite{nist2024} që standardizon ML-KEM (Module-Lattice-Based Key-Encapsulation Mechanism), i njohur më parë si CRYSTALS-Kyber. Ky algoritëm bazohet në problemin e vështirë të "Module Learning with Errors" (MLWE) dhe ofron tre nivele sigurie:

\begin{itemize}
\item \textbf{ML-KEM-512}: NIST Level 1 (sigurinë 128-bit)
\item \textbf{ML-KEM-768}: NIST Level 3 (sigurinë 192-bit)
\item \textbf{ML-KEM-1024}: NIST Level 5 (sigurinë 256-bit)
\end{itemize}

Qëllimi kryesor i këtij punimi është të demonstrojë se \textbf{kompjuterët klasikë mund të implementojnë dhe përdorin algoritme post-kuantike për të mbrojtur të dhënat nga kërcënimet e kompjuterëve kuantikë}, duke ofruar një zgjidhje praktike dhe të gatshme për përdorim.

\section{Baza Teorike}

\subsection{Kërcënimi Kuantik}

Kompjuterët kuantikë përdorin parimin e superpozicionit kuantik për të procesuar informacion në mënyrë paralele. Algoritmi i Shor mund të faktorizojë numra të mëdhenj primar dhe të zgjidhë problemin e logaritmit diskret në kohë polinomiale $O((\log N)^3)$, duke e bërë të mundur thyerjen e:

\begin{itemize}
\item \textbf{RSA}: Bazohet në vështirësinë e faktorizimit
\item \textbf{ECDH/ECDSA}: Bazohen në problemin e logaritmit diskret në kurba eliptike
\item \textbf{DH}: Bazohet në logaritmin diskret
\end{itemize}

Tabela \ref{tab:breaktime} paraqet krahasimin e kohës së nevojshme për thyerje:

\begin{table}[htbp]
\caption{Koha e Thyerjes së Algoritmeve}
\label{tab:breaktime}
\begin{center}
\begin{tabular}{|l|c|c|}
\hline
\textbf{Algoritmi} & \textbf{Kompjuter Klasik} & \textbf{Kompjuter Kuantik} \\
\hline
RSA-2048 & $10^9$ vjet & 8 orë* \\
ECC-256 & $10^{12}$ vjet & 10 minuta* \\
AES-256 & $2^{256}$ operacione & $2^{128}$ operacione \\
ML-KEM-768 & I pathyeshëm & I pathyeshëm \\
\hline
\multicolumn{3}{l}{\footnotesize *Me kompjuter kuantik me $\sim$4000 qubit të qëndrueshëm}
\end{tabular}
\end{center}
\end{table}

\subsection{Kriptografia Post-Kuantike}

Kriptografia post-kuantike përfshin algoritme që janë:
\begin{enumerate}
\item \textbf{Rezistente ndaj sulmeve kuantike}: Bazohen në probleme që nuk mund të zgjidhen efikasshëm as nga kompjuterët kuantikë
\item \textbf{Ekzekutueshme në kompjuterë klasikë}: Nuk kërkojnë hardware kuantik
\item \textbf{Praktike për përdorim}: Ofrojnë performancë të pranueshme
\end{enumerate}

Familjet kryesore të algoritmeve post-kuantike janë:
\begin{itemize}
\item \textbf{Lattice-based}: ML-KEM, ML-DSA (bazuar në rrjeta matematike)
\item \textbf{Hash-based}: SLH-DSA (bazuar në funksione hash)
\item \textbf{Code-based}: McEliece (bazuar në kode të korrigjimit të gabimeve)
\end{itemize}

\subsection{ML-KEM (CRYSTALS-Kyber)}

ML-KEM është një mekanizëm i enkapsulimit të çelësit (KEM) që mundëson shkëmbimin e sigurt të çelësave simetrikë \cite{bos2018}. Ai bazohet në problemin MLWE:

\textbf{Problemi MLWE}: Dallimi midis $(A, A \cdot s + e)$ dhe $(A, u)$, ku:
\begin{itemize}
\item $A$ është një matricë e rastësishme
\item $s$ është vektori sekret
\item $e$ është vektori i gabimit (noise)
\item $u$ është vektor i rastësishëm
\end{itemize}

Ky problem konsiderohet i vështirë edhe për kompjuterët kuantikë, duke ofruar siguri të qëndrueshme afatgjatë.

\textbf{Operacionet kryesore të ML-KEM:}
\begin{enumerate}
\item \textbf{KeyGen()}: Gjeneron çiftin $(pk, sk)$
\item \textbf{Encapsulate(pk)}: Krijon $(c, K)$ - ciphertext dhe çelësin e përbashkët
\item \textbf{Decapsulate(sk, c)}: Rikupton çelësin $K$
\end{enumerate}

\section{Arkitektura e Sistemit}

\subsection{Përmbledhje e Sistemit}

Sistemi i implementuar përbëhet nga tre komponentë kryesorë:

\begin{enumerate}
\item \textbf{Frontend (Angular 18)}: Ndërfaqja e përdoruesit për gjenerimin e çelësave, enkriptim dhe dekriptim
\item \textbf{Backend (.NET 9 API)}: Servisi kriptografik që implementon ML-KEM
\item \textbf{Biblioteka Kriptografike (BouncyCastle)}: Implementimi i algoritmit ML-KEM
\end{enumerate}

Fig. \ref{fig:architecture} paraqet arkitekturën e sistemit.

\begin{figure}[htbp]
\centerline{
\begin{tabular}{|c|}
\hline
\textbf{FRONTEND (Angular 18)} \\
Keys | Encrypt | Decrypt | Info \\
\hline
$\downarrow$ HTTP/REST \\
\hline
\textbf{BACKEND (.NET 9 API)} \\
PostQuantumCryptoController \\
PostQuantumCryptoService \\
\hline
$\downarrow$ \\
\hline
\textbf{BouncyCastle} \\
ML-KEM-512/768/1024 \\
\hline
\end{tabular}
}
\caption{Arkitektura e sistemit ML-KEM}
\label{fig:architecture}
\end{figure}

\subsection{Enkriptimi Hibrid}

Sistemi përdor qasjen hibride ku ML-KEM shkëmben çelësin simetrik, ndërsa AES-GCM enkripton të dhënat:

\begin{lstlisting}[language=C,caption={Enkriptimi Hibrid ML-KEM + AES-GCM}]
// Enkapsulo sekretin me ML-KEM
var enc = new MLKemEncapsulator(parameters);
enc.Init(pub);
enc.Encapsulate(kemCipher, shared);

// Enkripto me AES-GCM
var nonce = RandomNumberGenerator.GetBytes(12);
using (var aes = new AesGcm(shared, 16))
    aes.Encrypt(nonce, pt, ct, tag);
\end{lstlisting}

\subsection{Nivelet e Sigurisë}

Tabela \ref{tab:parameters} paraqet parametrat e tre niveleve të ML-KEM:

\begin{table}[htbp]
\caption{Parametrat e ML-KEM}
\label{tab:parameters}
\begin{center}
\begin{tabular}{|l|c|c|c|}
\hline
\textbf{Parametri} & \textbf{512} & \textbf{768} & \textbf{1024} \\
\hline
Niveli NIST & Level 1 & Level 3 & Level 5 \\
Siguria (bit) & 128 & 192 & 256 \\
Çelësi Publik & 800 B & 1184 B & 1568 B \\
Çelësi Privat & 1632 B & 2400 B & 3168 B \\
Ciphertext & 768 B & 1088 B & 1568 B \\
\hline
\end{tabular}
\end{center}
\end{table}

\section{Implementimi}

\subsection{Teknologjitë e Përdorura}

\textbf{Backend:}
\begin{itemize}
\item .NET 9.0 (ASP.NET Core Web API)
\item BouncyCastle.Cryptography 2.5.0
\item Swagger/OpenAPI për dokumentim
\end{itemize}

\textbf{Frontend:}
\begin{itemize}
\item Angular 18 (Standalone Components)
\item TypeScript 5.x
\item SCSS për stilizim
\end{itemize}

\subsection{API Endpoints}

Sistemi ekspozon katër endpoint-e REST:

\begin{verbatim}
GET  /api/PostQuantumCrypto/keys/{level}
POST /api/PostQuantumCrypto/encrypt
POST /api/PostQuantumCrypto/decrypt
GET  /api/PostQuantumCrypto/algorithms
\end{verbatim}

\section{Rezultatet dhe Analiza}

\subsection{Testimi Funksional}

Sistemi u testua me skenarë të ndryshëm të paraqitur në Tabelën \ref{tab:testing}:

\begin{table}[htbp]
\caption{Rezultatet e Testimit}
\label{tab:testing}
\begin{center}
\begin{tabular}{|l|c|}
\hline
\textbf{Skenari} & \textbf{Statusi} \\
\hline
Gjenerimi i çelësave (3 nivelet) & $\checkmark$ \\
Enkriptimi/Dekriptimi i tekstit & $\checkmark$ \\
Verifikimi i integritetit & $\checkmark$ \\
Trajtimi i gabimeve & $\checkmark$ \\
Testimi me mesazhe 1MB & $\checkmark$ \\
\hline
\end{tabular}
\end{center}
\end{table}

\subsection{Analiza e Performancës}

Tabela \ref{tab:performance} paraqet kohën e ekzekutimit:

\begin{table}[htbp]
\caption{Koha e Ekzekutimit (ms)}
\label{tab:performance}
\begin{center}
\begin{tabular}{|l|c|c|c|}
\hline
\textbf{Operacioni} & \textbf{512} & \textbf{768} & \textbf{1024} \\
\hline
KeyGen & 0.12 & 0.18 & 0.28 \\
Encapsulate & 0.15 & 0.22 & 0.35 \\
Decapsulate & 0.14 & 0.21 & 0.33 \\
Encrypt (1KB) & 0.25 & 0.32 & 0.45 \\
Decrypt (1KB) & 0.24 & 0.31 & 0.43 \\
\hline
\end{tabular}
\end{center}
\end{table}

\subsection{Krahasimi me RSA}

Tabela \ref{tab:comparison} tregon avantazhet e ML-KEM:

\begin{table}[htbp]
\caption{Krahasimi ML-KEM vs RSA}
\label{tab:comparison}
\begin{center}
\begin{tabular}{|l|c|c|c|}
\hline
\textbf{Metrika} & \textbf{RSA-2048} & \textbf{ML-KEM-768} & \textbf{Përparësia} \\
\hline
Siguria kuantike & Jo & Po & ML-KEM \\
Çelësi publik & 256 B & 1184 B & RSA \\
Koha gjenerimit & 150 ms & 0.18 ms & ML-KEM \\
Koha enkriptimit & 1 ms & 0.22 ms & ML-KEM \\
\hline
\end{tabular}
\end{center}
\end{table}

\section{Diskutimi}

\subsection{Avantazhet e Qasjes}

Implementimi demonstron disa avantazhe kyçe:

\begin{enumerate}
\item \textbf{Praktikiteti}: Sistemi funksionon plotësisht në kompjuterë klasikë pa nevojë për hardware të specializuar.

\item \textbf{Gatishmëria}: Organizatat mund të fillojnë migrimin në algoritme post-kuantike tani, para se kompjuterët kuantikë të bëhen kërcënim real.

\item \textbf{Pajtueshmëria}: Enkriptimi hibrid (ML-KEM + AES-GCM) siguron pajtueshmëri me infrastrukturën ekzistuese.

\item \textbf{Skalueshmëria}: Tre nivele sigurie mundësojnë balancimin e sigurisë dhe performancës sipas nevojave.
\end{enumerate}

\subsection{Sfidat dhe Kufizimet}

\begin{enumerate}
\item \textbf{Madhësia e çelësave}: Çelësat e ML-KEM janë më të mëdhenj se ata të RSA.
\item \textbf{Migrimi}: Sistemet legacy kërkojnë përditësim.
\item \textbf{Standardizimi}: Algoritmet shtesë ende po standardizohen.
\end{enumerate}

\subsection{Rekomandimet}

\begin{enumerate}
\item \textbf{Filloni tani}: "Harvest now, decrypt later" - të dhënat e sotme mund të dekriptohen nesër.
\item \textbf{Qasja hibride}: Kombinoni algoritmet post-kuantike me ato klasike.
\item \textbf{Planifikimi}: Identifikoni sistemet kritike dhe planifikoni migrimin gradual.
\end{enumerate}

\section{Përfundime}

Ky punim demonstroi sukseshshëm se kompjuterët klasikë mund të implementojnë dhe përdorin algoritme post-kuantike për të siguruar komunikime ndaj kërcënimeve të ardhshme kuantike. Implementimi i ML-KEM sipas NIST FIPS 203 ofron:

\begin{itemize}
\item \textbf{Sigurinë afatgjatë} ndaj kompjuterëve kuantikë
\item \textbf{Performancë të shkëlqyeshme} të krahasueshme me algoritmet klasike
\item \textbf{Praktikitet} për përdorim në sisteme reale
\end{itemize}

Sistemi i zhvilluar provon se tranzicioni në kriptografi post-kuantike nuk kërkon hardware kuantik - kompjuterët e sotëm janë plotësisht të aftë të ekzekutojnë këto algoritme me efikasitet. Organizatat duhet të fillojnë planifikimin e migrimit tani për të siguruar mbrojtje afatgjatë të të dhënave sensitive.

\textbf{Kontributi kryesor}: Demonstrimi praktik se sigurim post-kuantik është i arritshëm me teknologji aktuale, duke eliminuar barrierën e perceptuar se nevojitet hardware i specializuar.

\begin{thebibliography}{10}

\bibitem{shor1994}
P.~W. Shor, ``Algorithms for quantum computation: discrete logarithms and factoring,'' in \emph{Proc. 35th Annual Symposium on Foundations of Computer Science}, IEEE, 1994, pp. 124--134.

\bibitem{nas2019}
National Academies of Sciences, Engineering, and Medicine, \emph{Quantum Computing: Progress and Prospects}. Washington, DC: The National Academies Press, 2019.

\bibitem{nist2024}
National Institute of Standards and Technology, ``FIPS 203: Module-Lattice-Based Key-Encapsulation Mechanism Standard,'' August 2024. [Online]. Available: https://csrc.nist.gov/pubs/fips/203/final

\bibitem{bos2018}
J.~Bos \emph{et al.}, ``CRYSTALS-Kyber: A CCA-Secure Module-Lattice-Based KEM,'' in \emph{2018 IEEE European Symposium on Security and Privacy (EuroS\&P)}, 2018, pp. 353--367.

\bibitem{moody2022}
D.~Moody \emph{et al.}, ``Status Report on the Third Round of the NIST Post-Quantum Cryptography Standardization Process,'' NIST IR 8413, 2022.

\bibitem{rfc8446}
E.~Rescorla, ``The Transport Layer Security (TLS) Protocol Version 1.3,'' RFC 8446, IETF, August 2018.

\bibitem{nist800-38d}
M.~Dworkin, ``Recommendation for Block Cipher Modes of Operation: Galois/Counter Mode (GCM) and GMAC,'' NIST SP 800-38D, November 2007.

\bibitem{bouncycastle}
The Legion of the Bouncy Castle, ``BouncyCastle Cryptography Library Documentation,'' 2024. [Online]. Available: https://www.bouncycastle.org/

\bibitem{bernstein2017}
D.~J. Bernstein and T.~Lange, ``Post-quantum cryptography,'' \emph{Nature}, vol. 549, no. 7671, pp. 188--194, September 2017.

\bibitem{cisa2024}
CISA, NSA, and NIST, ``Quantum-Readiness: Migration to Post-Quantum Cryptography,'' Cybersecurity Information Sheet, August 2024.

\end{thebibliography}

\end{document}

